% =============================================================================
%                            Chapter 1: Introduction
% =============================================================================

\setcounter{chapter}{1}
\setcounter{figure}{0}
\setcounter{table}{0}
\chapter*{Chapter 1}
\markboth{Introduction}{Introduction} 
\section{Introduction}

Agriculture in Pakistan faces challenges like water shortage, low 
productivity, and lack of modern tools \cite{shafi2019precision}. 
To address this, the \textbf{Smart AI Powered Agriculture System} 
uses Internet of Things (IoT) sensors to monitor soil and weather 
conditions, automates irrigation to save water, and applies 
artificial intelligence (AI) to recommend suitable crops 
\cite{farooq2019survey}. This solution is designed to help farmers 
improve yield, reduce costs, save water, and adopt smarter farming 
practices.

\subsection{Project Introduction}

The Smart AI Powered Agriculture System is an IoT-based solution 
designed to make farming more efficient and sustainable. The system 
uses sensors to monitor soil moisture, temperature, humidity, light, 
and rainfall in real time. This data is sent to a mobile application 
through a backend server, allowing farmers to track field conditions, 
receive alerts, and view historical records \cite{mohanraj2016field}.

The system also includes automatic irrigation, where water is 
supplied only when the soil is dry \cite{garcia2020iot}, and an AI 
module that recommends the most suitable crops for upcoming seasons 
\cite{musanase2023data}. The goal is to help farmers conserve water, 
improve crop yield, and make data-driven decisions in agriculture.

The main beneficiaries of this project are small and medium scale 
farmers in Pakistan, who often face challenges such as water wastage, 
low crop productivity, and lack of access to modern agricultural 
tools \cite{zia2021experimental}. By using this system, farmers can 
monitor their fields easily, save water through automated irrigation, 
and receive crop recommendations that support better planning.

\subsection{Existing Examples / Solutions}

Several existing systems and solutions in Pakistan and globally work 
towards precision agriculture, smart irrigation, and farm monitoring. 
However, most of them are either expensive, focused on large-scale 
farms, or limited in scope \cite{khanna2019evolution}. 
Table~\ref{tab:existing_examples} summarizes some examples and the 
gaps that remain.

\begin{table}[H]
\centering
\caption{Existing Examples}
\label{tab:existing_examples}
\begin{tabular}{|p{3.2cm}|p{5.2cm}|p{5.2cm}|}
\hline
\textbf{Name} & \textbf{What they do (features)} & 
\textbf{What gap remains} \\ \hline

SAWiE (Sustainable Agriculture Water \& Intelligent Ecosystem) &
Offers IoT and machine learning based advisory for farmers; soil 
mapping and alerts for droughts and floods. &
Does not directly automate irrigation hardware or pump control. 
\\ \hline

Buraq Integrated Solutions --- Smart Drip Irrigation &
Implements drip irrigation systems with soil moisture and temperature 
sensors. &
Lacks an AI-based crop recommendation module. \\ \hline

VGreen --- CropSight\textsuperscript{\texttrademark} &
Provides farm mapping, real-time monitoring, and analytics; localized 
for Pakistani farms. &
Good dashboards and analytics but lacks full automation (pump control 
and crop recommendation) and is costly for small farmers. \\ \hline

Field Commander --- Valley Irrigation Pakistan &
Offers remote monitoring and SMS/Email alerts. &
Does not offer smart irrigation and is mainly available for large 
commercial farms. \\ \hline

Smart IoT Farm at PMAS--AAUR, Rawalpindi &
University research farm with sensors (e.g., soil moisture) and 
real-time IoT monitoring to demonstrate precision agriculture. &
No dedicated mobile application for farmers, no AI-based decision 
support, and no automated irrigation. \\ \hline
\end{tabular}
\end{table}

The Smart AI Powered Agriculture System aims to combine IoT-based 
sensing, automatic irrigation, mobile app support, and AI-driven 
crop recommendations in a single affordable solution targeted at 
small and medium farmers.

\subsection{Problem Statement}

Agriculture is the main source of income and food in many countries, 
especially in developing regions. Many farmers depend on farming for 
their livelihood, but they face challenges such as water wastage, 
unpredictable weather, limited access to real-time information, and 
lack of proper guidance for choosing crops or managing their fields 
\cite{shafi2019precision}. Traditional farming methods often result 
in over-watering or under-watering, loss of soil nutrients, and low 
crop production \cite{garcia2020iot}. Even though modern technology 
is available, many small and medium farmers cannot afford or access 
smart farming tools. There is a need for an affordable and 
easy-to-use smart agriculture system that uses IoT and AI to monitor 
the environment, save resources such as water and energy, and provide 
useful recommendations to farmers.

\subsection{Business Scope}

The Smart AI Powered Agriculture System has strong potential in the 
agriculture sector of Pakistan. Since many farmers face problems like 
water wastage, low crop yield, and lack of modern tools, this system 
provides an affordable and simple solution.

By offering IoT-based smart farming kits (sensors, controller, and 
mobile application), farmers can save water, achieve better crop 
production, and make smarter decisions. The system can be adopted by 
small and medium farmers and can also be promoted through government 
programs, non-governmental organizations (NGOs), and agriculture 
companies. The idea has strong business value because it directly 
improves farmers' income and supports food security.

\begingroup \footnotesize \renewcommand{\arraystretch}{1.15} \begin{longtable}{|p{3.8cm}|p{3.8cm}|p{3.8cm}|} \caption{Lean Canvas} \label{tab:lean_canvas} \\ \hline \textbf{Problem} & \textbf{Solution} & \textbf{Unique Value Proposition} \\ \hline \endfirsthead \hline \endhead \hline \multicolumn{3}{r}{\small\textit{Continued on next page}} \\ \endfoot \hline \endlastfoot 1. Water wastage due to over-irrigation. \newline 2. Low crop yield from poor crop choice. \newline 3. Lack of affordable smart farming tools. & 1. IoT sensors for soil and weather. \newline 2. Automatic irrigation system. \newline 3. Mobile app with AI-based crop suggestion. & An affordable smart farming kit that saves water, increases yield, and helps farmers make better crop decisions. \\ \hline \textbf{Existing Alternatives} & \textbf{Key Metrics (Approx.)} & \textbf{High-Level Concept} \\ \hline 1. Manual irrigation. \newline 2. Expensive imported systems. & 1. 100+ farmers use the system. \newline 2. 45\% water savings (est.). \newline 3. 35\% increase in yield (est.). & A smart farming assistant for every farmer, like a personal digital guide for crops. \\ \hline \textbf{Channels} & \textbf{Cost Structure} & \textbf{Customer Segments} \\ \hline 1. Direct sales to farmers. \newline 2. Partnerships with NGOs and government programs. \newline 3. Agriculture supply companies. & 1. IoT hardware (sensors, controllers). \newline 2. Mobile app and server hosting. \newline 3. Installation and support. & 1. Small and medium farmers. \newline 2. Agriculture NGOs and government projects. \\ \hline \textbf{Early Adopters} & \multicolumn{2}{p{\dimexpr7.6cm+2\tabcolsep+\arrayrulewidth}|} {\textbf{Revenue Structure}} \\ \hline 1. Farmers with 1--4 acres of land. \newline 2. NGOs running pilot agriculture projects. & \multicolumn{2}{p{\dimexpr7.6cm+2\tabcolsep+\arrayrulewidth}|}{ 1. Selling hardware kits. \newline 2. Subscription for premium app features. \newline 3. Installation and training services. } \\ \hline \end{longtable} \endgroup

\subsection{Useful Tools and Technologies}
The key hardware and software components for the Smart AI Powered 
Agriculture System are:

\subsubsection{Hardware}
The hardware architecture is anchored by the ESP32 microcontroller, which serves as the central processing unit equipped with built-in Wi-Fi capabilities to facilitate seamless sensor data processing and transmission \cite{raj2021survey}. To accurately monitor the environmental conditions, a capacitive soil moisture sensor (YL-69) is deployed to detect the volumetric water content in the soil. Additionally, a DHT22 sensor provides high-precision measurements of ambient temperature and humidity, while a Light Dependent Resistor (LDR) continuously monitors the sunlight intensity essential for optimal crop growth. To prevent over-watering during natural precipitation, a rain sensor (FC-37) is integrated into the system. Actuation is managed by a dedicated relay module that safely controls the higher-voltage water pump based on the microcontroller's logic.

\subsubsection{Software}
The software ecosystem spans multiple layers, beginning with the firmware developed using C++ in the Arduino IDE to program the ESP32 microcontroller. At the backend, a robust Node.js server is deployed to efficiently handle incoming telemetry data from the field sensors and manage bi-directional communication with the mobile application. All sensor readings, historical records, and user configurations are persistently stored in a normalized MySQL relational database. For the frontend user experience, a cross-platform mobile application is developed using the Flutter framework (or native Android Java/Kotlin), which presents live data, historical trends, and system alerts through an intuitive interface. To enhance the analytical value of the data, visualization libraries (Chart.js) is utilized to render dynamic graphs and charts. Furthermore, real-time push notifications regarding critical field conditions are delivered to the farmers via Firebase Cloud Messaging \cite{khawas2018firebase}.

\subsubsection{AI and Machine Learning}
The predictive intelligence of the system is driven by a machine learning module developed in Python's Random Forest Classifier model to generate accurate crop recommendations and conduct predictive analytics based on historical environmental data \cite{musanase2023data}. To ensure the high fidelity of the inputs fed into the models, advanced signal processing techniques are applied to filter out noise and clean the raw sensor data prior to analysis.

\subsection{Project Work Break Down}
The project is divided into multiple phases to ensure systematic 
development, testing, and deployment. Rather than treating the system 
as a single monolithic build, the Work Breakdown Structure (WBS) 
decouples the architecture into discrete, parallelizable engineering 
tracks \cite{sommerville2016software}. 

The core tracks include: (1) Hardware and Firmware, encompassing the 
ESP32 sensor integration and the 3-tier polling safety logic; (2) 
Backend Infrastructure, focusing on the Node.js API and MySQL 
relational database schema; (3) Artificial Intelligence, detailing the 
synthesis of the Gaussian dataset and the training of the Random Forest 
classifier; and (4) Frontend Application, covering the Flutter mobile 
interface for live dashboard monitoring and manual irrigation control.


To address the detailed task distribution, Table~\ref{tab:wbs_detailed} outlines the formal Work Breakdown Structure with specific task ownership across the team members (Muhammad Awais, Hamza Bashir, Junaid Amin), explicit phase durations, and deliverables.

\renewcommand{\arraystretch}{1.4}

\begin{longtable}{|p{0.8cm}|p{3.5cm}|p{2cm}|p{3.5cm}|p{4.5cm}|}
\caption{Detailed Work Breakdown Structure with Ownership}
\label{tab:wbs_detailed}\\

\hline
\rowcolor{black}
\color{white}\textbf{ID} &
\color{white}\textbf{Task / Phase} &
\color{white}\textbf{Duration} &
\color{white}\textbf{Ownership} &
\color{white}\textbf{Deliverable} \\
\hline
\endfirsthead

\multicolumn{5}{c}%
{\tablename\ \thetable\ -- Continued from previous page} \\

\hline
\rowcolor{black}
\color{white}\textbf{ID} &
\color{white}\textbf{Task / Phase} &
\color{white}\textbf{Duration} &
\color{white}\textbf{Ownership} &
\color{white}\textbf{Deliverable} \\
\hline
\endhead

\hline
\multicolumn{5}{r}{Continued on next page} \\
\hline
\endfoot

\hline
\endlastfoot

\textbf{1} & \textbf{Hardware \& Firmware} & \textbf{4 Weeks} & \textbf{Hamza Bashir} & \textbf{Physical prototype \& ESP32 Code} \\ \hline
1.1 & Sensor Calibration & 1 Week & Hamza Bashir & Calibrated DHT22 and Moisture inputs \\ \hline
1.2 & Actuator Integration & 1 Week & Hamza Bashir & Functional relay control for pump \\ \hline
1.3 & ESP32 Firmware & 2 Weeks & Hamza Bashir & Stable Wi-Fi telemetry and control loop \\ \hline

\textbf{2} & \textbf{Backend \& Database} & \textbf{4 Weeks} & \textbf{Junaid Amin} & \textbf{Deployed API \& MySQL Schema} \\ \hline
2.1 & Database Design & 1 Week & Junaid Amin & Normalized schema deployed to server \\ \hline
2.2 & REST API Dev & 2 Weeks & Junaid Amin & Endpoints for telemetry and control \\ \hline
2.3 & Authentication & 1 Week & Junaid Amin & Secure JWT login and field mapping \\ \hline

\textbf{3} & \textbf{AI \& Machine Learning} & \textbf{3 Weeks} & \textbf{Muhammad Awais} & \textbf{Trained Scikit-Learn Model} \\ \hline
3.1 & Dataset Preparation & 1 Week & Muhammad Awais & Dataset generation and cleaning \\ \hline
3.2 & Model Training & 1 Week & Muhammad Awais & Tuned Random Forest (n\_estimators=150) \\ \hline
3.3 & API Integration & 1 Week & Muhammad Awais & Wrapper returning confidence score \\ \hline

\textbf{4} & \textbf{Frontend Application} & \textbf{4 Weeks} & \textbf{Team} & \textbf{Compiled Flutter APK} \\ \hline
4.1 & UI Mockups \& Flow & 1 Week & Junaid Amin & Figma designs for all screens \\ \hline
4.2 & Core Development & 2 Weeks & Team & Flutter UI and backend API binding \\ \hline
4.3 & System Testing & 1 Week & Hamza Bashir & Passed test cases across devices \\ \hline

\end{longtable}

\subsection{Project Time Line}
The project timeline spans from 29 September 2025 to 10 July 2026. 
The schedule allocates dedicated time blocks for isolated component 
development followed by rigorous integration phases. For instance, 
the AI model training phase occurs concurrently with backend API 
development, ensuring that the prediction endpoints are immediately 
consumable once the model achieves its target accuracy metrics. 
Similarly, firmware development prioritizes local offline safety 
logic before network connectivity is fully integrated, reducing 
critical path dependencies.

To provide a precise technical schedule, Table~\ref{tab:project_schedule} breaks down the project timeline into formal milestones, mapping durations to exact dates, defining cross-component dependencies, and establishing strict team ownership.

\enlargethispage{5cm}
\begin{table}[H]
    \centering
    \caption{Formal Project Schedule}
    \label{tab:project_schedule}
    \vspace{0.2cm}
    \renewcommand{\arraystretch}{1.4}
    \begin{tabularx}{\textwidth}{|c|X|c|c|c|c|}
        \hline
        \rowcolor{black}
        \color{white}\textbf{ID} & \color{white}\textbf{Milestone} & \color{white}\textbf{Start} & \color{white}\textbf{End} & \color{white}\textbf{Depends On} & \color{white}\textbf{Owner} \\ \hline
        \textbf{M1}  & Requirement Specifications & 29 Sep 25 & 12 Oct 25 & -- & Team \\ \hline
        \textbf{M2}  & Hardware Design & 13 Oct 25 & 02 Nov 25 & M1 & Hamza \\ \hline
        \textbf{M3}  & Backend \& Database & 03 Nov 25 & 30 Nov 25 & M1 & Junaid \\ \hline
        \textbf{M4}  & AI Dataset Preparation & 13 Oct 25 & 26 Oct 25 & M1 & Awais \\ \hline
        \textbf{M5}  & ESP32 Firmware & 03 Nov 25 & 23 Nov 25 & M2 & Hamza \\ \hline
        \textbf{M6}  & AI Model Training & 27 Oct 25 & 16 Nov 25 & M4 & Awais \\ \hline
        \textbf{M7}  & API Integration & 01 Dec 25 & 21 Dec 25 & M3, M5, M6 & Team \\ \hline
        \textbf{M8}  & Mobile App UI & 22 Dec 25 & 18 Jan 26 & M1 & Junaid \\ \hline
        \textbf{M9}  & System Integration & 19 Jan 26 & 15 Mar 26 & M7, M8 & Team \\ \hline
        \textbf{M10} & Testing \& Fixing & 16 Mar 26 & 10 May 26 & M9 & Hamza \\ \hline
        \textbf{M11} & Final Deployment & 11 May 26 & 10 Jul 26 & M10 & Team \\ \hline
    \end{tabularx}
\end{table}

\setcounter{chapter}{1}